% CHOICE (pref-1): series naming-preface prepended as front matter; the Vol-2 "Oracle of Free Will"
% preface (00-preface.tex) is retained intact immediately after it (the Ch30 ending depends on it).
% CHOICE (pref-2): this preface NAMES the question --- choice, the selector, re-representation; the
% historical spine Zermelo -> Frege -> Wittgenstein -> Cohen; and "the axiom of choice" itself (the
% Axiom of Choice section). It WITHHOLDS the resolution the body owns: it does not identify the
% instrument's method with Cohen forcing, does not assert choice-freeness as the answer, and names none
% of the Vol-2 reserved reveals (the fact->verdict turn, the sign a turned loop carries, the spinor, the
% continuum floor, the verdict). The preface poses; the body, at its climax (Ch29 and Ch30), resolves.

\chapter*{Naming}
\addcontentsline{toc}{chapter}{Series Preface --- Naming}
\markboth{Series Preface}{Series Preface}

\begin{center}\textit{Series Preface}\end{center}

\bigskip

\section*{The Instrument}

The books in this series build one instrument, and build it the same way each time. It takes the world only as
the world arrives: in finite marks, a count here, a coincidence there --- nothing continuous ever assumed and
nothing unbounded ever required. It carries no endless tape and no perfect ruler. Whatever it cannot tell apart
at the resolution it actually has, it treats as one, for it has no finer place to stand. From that scant
equipment it reads what can be honestly read, and it declines, on principle, to read anything else.

Because it is finite, it never reaches inside itself for an answer the world did not hand it. It does not guess;
it does not pick a representative out of the air; it does not help itself to a choice it cannot account for.
Everything it reports, it builds from what is already on the record. This modesty is not a limitation the series
apologizes for --- it is the discipline the series is about. An instrument that will not choose freely turns out
to have a great deal to say about choice.

\section*{The Demon}

There is an older instrument this one is measured against. Laplace imagined a mind that held, at a single
instant, the place and motion of every particle there is and every force between them, and could read from that
one state the whole of the future and the whole of the past. To such a mind the world is a function: hand it the
present, and it returns the next moment, and the next, without remainder and without choice. Nothing is open,
because nothing was ever unfixed.

The instrument the series builds cannot be that mind, and its inability is exact rather than vague. It holds the
present only to a finite sharpness, and beneath that floor two states it cannot separate run off into futures it
also cannot separate. And where the demon returned a single next moment, the instrument returns a family of them
--- the admissible ones, the ones the record does not forbid --- and waits for the world to do the selecting. A
function where the demon stood; a family, and an act, where the instrument stands. The distance between the two
is the width of everything this series has to say.

\section*{The Axiom of Choice}

Name that distance plainly, at its origin. To name a thing is, first, to pick a representative --- to let one
member of a collection stand for the whole of it. Early in the last century Zermelo raised the picking to a
principle: that from any family of nonempty collections one may always draw a representative, one from each, all
at once. Two things are true of the drawn representative. It \emph{exists} --- the principle guarantees it. And
it is chosen \emph{freely} --- nothing inside any collection says which of its members the choice must fall on.
This is the axiom of choice, and it is the quiet assumption beneath a great deal of mathematics and, if one is
not careful, beneath a great deal of physics. It asserts a selector without saying how the selector selects: a
hand that reaches into each collection and comes out holding one thing, guided by nothing on the record.

And there is more to a name than the thing it lands on. Frege saw that a name carries also a \emph{sense} --- a
way the thing is presented, a manner of being picked out. The morning star and the evening star are one planet
under two presentations; to name is not merely to seize an element but to seize it under some mode. So the choice
the axiom asks for is wider than it first appears: not only which member, but under which presentation the member
is taken. The series takes this widened choice as its central question --- not to deny the axiom, for as
mathematics it stands, but to ask what becomes of a physics that will not pay for it: what an instrument can
still say when it refuses to reach into the world and pick.

\section*{Physics as Selection}

The refusal turns out to be productive, because the world does the reaching for you. Again and again the
instrument runs its account to the end, settles everything the record fixes, and arrives at a last step it cannot
take from the inside --- a representative it would have to choose, and cannot. At exactly that step the world
supplies what the instrument withheld. Which outcome, of the admissible family; which member, of the collection
the account left open --- the world selects, at the moment of measurement, and the instrument reads the selection
off rather than making it.

So the missing selector is not missing after all; it was never the instrument's to provide. Physics, on this
telling, is the study of a selection the world performs and the record receives. The axiom asked for a chooser
and named none. The instrument names the chooser: it is the world, met at the point of measurement, and never
the desk.

\section*{The Locality of Choice}

And the selection is local. It is not read off a global coordinate laid down once over everything; there is no
such coordinate, and an instrument that went looking for one would find nothing there. The selection is real only
where it is made --- at the particular reading, in the particular act, answerable to the particular fact in front
of it. Meaning, the later Wittgenstein taught, is use: a name means what the practice does with it, not what a
private hand assigns it in advance. The instrument's selections are of that kind. There is no fact of the matter
about the choice apart from the making of it; the choice is what happens at the site, and nowhere else, and it
cannot be carried off somewhere and settled in advance.

And some choices the accepted rules can neither compel nor forbid. Cohen showed, at the very foundations, that
certain questions the standing axioms leave genuinely open --- a place where an answer is expected and none is
forced, and where either answer could be built outward from what is already fixed, without contradiction. The
choice of representation, at such a place, is not decreed by the rules; it is approached by construction, raised
from the public conditions already on the record. What the instrument does at exactly such a step --- how it
reaches an answer the rules do not hand it, out of finite public structure and nothing more --- is the method
these volumes exist to display, and a preface has no business anticipating it. Between Zermelo's free reach and
this local, public making runs a change the series will spend its length earning: the choice that began as a
private pick becomes an act performed in the open, at a place, answerable to a record --- the widest freedom
narrowed to the one the situation compels. Between the naming of that question and its answering there is a
distance, and the distance is the book.

% CHOICE (pref-2): the draft named this section "Free Will"; renamed to "The Last Facing" so the series
% preface does NOT deliver the words "free will" before the Vol-2 "Oracle of Free Will" preface's
% load-bearing line ("I will set it down once, here ... free will"). The section poses the theme and hands
% the naming to the volume; the crisp term stays the Oracle preface's to set down first.
\section*{The Last Facing}

There is one selection the series measures everything around and does not itself make. When the instrument has
run its accounts to the end and the world has supplied every selection the record left open, one facing remains
that no measurement fixes and no law compels --- the direction in which a reader turns to meet what has been
counted. The count is the world's, complete and handed over; the meeting is the reader's. Each volume arrives at
this in its own physics and does the same scrupulous thing: it settles everything it can, and leaves that last
facing to the one who is reading. What to call it, and how much weight it can bear, the volume before you will
name in its own turn, and at its end hand back. The series insists only on the shape --- a reckoning that earns
every step but the last, and declines, on purpose, to take it.

\section*{What Not to Claim}

An honest instrument names the floor it stands on and reports nothing beneath it, and a series built on such an
instrument owes its reader that same honesty at the outset. So, plainly, what these books do not claim. They do
not claim the world is finite because the instrument is; the finiteness is the instrument's discipline, not a
metaphysics smuggled in. They do not claim to have derived, from physics, the freedom named a moment ago, nor to
have proved it, nor to have proved its absence; the instrument reaches a floor and stops, and does not pretend to
see beneath it. They do not claim that where the record leaves a step open, the world fills it at random, or by
whim, or by anything the books can themselves characterize; they claim only that the record does not fill it, and
that the filling, when it comes, comes from outside the account. And they claim no authority over what a reader
makes of the reading. The instrument reports what is. What to do about it was never an instrument's to say.

\section*{The Final Turn}

Everything the series measures, it measures on its way to one last step it will not take. The accounts come in;
the selections the world makes are read off, local and forced and complete; and at the end there stands a single
act of selection that no prior selection settles --- the choosing of how to face a world that has been, in every
other respect, fully counted. That act the books hand over unmade. It is the one representative the instrument,
having refused all the others, still refuses to pick on anyone's behalf --- for to pick it would be to become the
demon after all, and the whole series is built to show that the demon cannot be built. So the last selection is
left where it belongs. The world can be named. The naming of what to do with the world is the reader's, and it is
the one naming these books decline to perform.

\bigskip

\noindent What follows, in this volume, is a physics book. It begins again from the beginning --- from a single
mark, and the one line drawn through all of them, between what the world fixes and what is ours to set --- and it
does not mention the question named here until it has earned the right to. Read once, each volume is an account
of the world. Read twice, it is the story of a name that would not be chosen, and had to be constructed instead.
